% ============================================================
% fig_perceptron.tex: Diagrama del perceptrón
% Uso: % ============================================================
% fig_perceptron.tex: Diagrama del perceptrón
% Uso: % ============================================================
% fig_perceptron.tex: Diagrama del perceptrón
% Uso: % ============================================================
% fig_perceptron.tex: Diagrama del perceptrón
% Uso: \input{figures/fig_perceptron}
% ============================================================
\begin{figure}[H]
\centering
\begin{tikzpicture}[
    node distance=1.4cm,
    input/.style={circle, draw=gray!60, fill=gray!15, minimum size=0.8cm, font=\small},
    neuron/.style={circle, draw=blue!60, fill=blue!10, minimum size=1.0cm, font=\small\bfseries},
    output/.style={circle, draw=green!60!black, fill=green!10, minimum size=0.8cm, font=\small},
    arrow/.style={-{Stealth}, thick}
]

% Entradas
\node[input] (x1) at (0, 2)   {$x_1$};
\node[input] (x2) at (0, 0.7) {$x_2$};
\node        (xd) at (0,-0.1) {$\vdots$};
\node[input] (xn) at (0,-1)   {$x_n$};

% Neurona
\node[neuron] (n) at (3, 0.7) {$\Sigma$};

% Salida
\node[output] (y) at (5.5, 0.7) {$y$};

% Bias
\node[input] (b) at (1.5, -1.8) {$b$};

% Función de activación (caja)
\node[draw=orange!70, fill=orange!10, rounded corners, minimum width=1.1cm,
      minimum height=0.7cm, font=\small] (phi) at (4.4, 0.7) {$\varphi$};

% Flechas entradas → neurona
\draw[arrow] (x1) -- node[above, font=\tiny] {$w_1$} (n);
\draw[arrow] (x2) -- node[above, font=\tiny] {$w_2$} (n);
\draw[arrow] (xn) -- node[below, font=\tiny] {$w_n$} (n);
\draw[arrow] (b)  -- (n);

% Neurona → activación → salida
\draw[arrow] (n)   -- (phi);
\draw[arrow] (phi) -- (y);

% Etiquetas
\node[font=\footnotesize, gray] at (0, -2.2) {Entradas};
\node[font=\footnotesize, gray] at (3, -2.2) {Suma ponderada};
\node[font=\footnotesize, gray] at (5.5, -2.2) {Salida};

\end{tikzpicture}
\caption{Estructura del perceptrón. Las entradas $x_i$ se combinan linealmente
con pesos $w_i$ y un sesgo $b$; la suma se pasa por una función de activación
$\varphi$ para producir la salida $y$.}
\label{fig:perceptron}
\end{figure}

% ============================================================
\begin{figure}[H]
\centering
\begin{tikzpicture}[
    node distance=1.4cm,
    input/.style={circle, draw=gray!60, fill=gray!15, minimum size=0.8cm, font=\small},
    neuron/.style={circle, draw=blue!60, fill=blue!10, minimum size=1.0cm, font=\small\bfseries},
    output/.style={circle, draw=green!60!black, fill=green!10, minimum size=0.8cm, font=\small},
    arrow/.style={-{Stealth}, thick}
]

% Entradas
\node[input] (x1) at (0, 2)   {$x_1$};
\node[input] (x2) at (0, 0.7) {$x_2$};
\node        (xd) at (0,-0.1) {$\vdots$};
\node[input] (xn) at (0,-1)   {$x_n$};

% Neurona
\node[neuron] (n) at (3, 0.7) {$\Sigma$};

% Salida
\node[output] (y) at (5.5, 0.7) {$y$};

% Bias
\node[input] (b) at (1.5, -1.8) {$b$};

% Función de activación (caja)
\node[draw=orange!70, fill=orange!10, rounded corners, minimum width=1.1cm,
      minimum height=0.7cm, font=\small] (phi) at (4.4, 0.7) {$\varphi$};

% Flechas entradas → neurona
\draw[arrow] (x1) -- node[above, font=\tiny] {$w_1$} (n);
\draw[arrow] (x2) -- node[above, font=\tiny] {$w_2$} (n);
\draw[arrow] (xn) -- node[below, font=\tiny] {$w_n$} (n);
\draw[arrow] (b)  -- (n);

% Neurona → activación → salida
\draw[arrow] (n)   -- (phi);
\draw[arrow] (phi) -- (y);

% Etiquetas
\node[font=\footnotesize, gray] at (0, -2.2) {Entradas};
\node[font=\footnotesize, gray] at (3, -2.2) {Suma ponderada};
\node[font=\footnotesize, gray] at (5.5, -2.2) {Salida};

\end{tikzpicture}
\caption{Estructura del perceptrón. Las entradas $x_i$ se combinan linealmente
con pesos $w_i$ y un sesgo $b$; la suma se pasa por una función de activación
$\varphi$ para producir la salida $y$.}
\label{fig:perceptron}
\end{figure}

% ============================================================
\begin{figure}[H]
\centering
\begin{tikzpicture}[
    node distance=1.4cm,
    input/.style={circle, draw=gray!60, fill=gray!15, minimum size=0.8cm, font=\small},
    neuron/.style={circle, draw=blue!60, fill=blue!10, minimum size=1.0cm, font=\small\bfseries},
    output/.style={circle, draw=green!60!black, fill=green!10, minimum size=0.8cm, font=\small},
    arrow/.style={-{Stealth}, thick}
]

% Entradas
\node[input] (x1) at (0, 2)   {$x_1$};
\node[input] (x2) at (0, 0.7) {$x_2$};
\node        (xd) at (0,-0.1) {$\vdots$};
\node[input] (xn) at (0,-1)   {$x_n$};

% Neurona
\node[neuron] (n) at (3, 0.7) {$\Sigma$};

% Salida
\node[output] (y) at (5.5, 0.7) {$y$};

% Bias
\node[input] (b) at (1.5, -1.8) {$b$};

% Función de activación (caja)
\node[draw=orange!70, fill=orange!10, rounded corners, minimum width=1.1cm,
      minimum height=0.7cm, font=\small] (phi) at (4.4, 0.7) {$\varphi$};

% Flechas entradas → neurona
\draw[arrow] (x1) -- node[above, font=\tiny] {$w_1$} (n);
\draw[arrow] (x2) -- node[above, font=\tiny] {$w_2$} (n);
\draw[arrow] (xn) -- node[below, font=\tiny] {$w_n$} (n);
\draw[arrow] (b)  -- (n);

% Neurona → activación → salida
\draw[arrow] (n)   -- (phi);
\draw[arrow] (phi) -- (y);

% Etiquetas
\node[font=\footnotesize, gray] at (0, -2.2) {Entradas};
\node[font=\footnotesize, gray] at (3, -2.2) {Suma ponderada};
\node[font=\footnotesize, gray] at (5.5, -2.2) {Salida};

\end{tikzpicture}
\caption{Estructura del perceptrón. Las entradas $x_i$ se combinan linealmente
con pesos $w_i$ y un sesgo $b$; la suma se pasa por una función de activación
$\varphi$ para producir la salida $y$.}
\label{fig:perceptron}
\end{figure}

% ============================================================
\begin{figure}[H]
\centering
\begin{tikzpicture}[
    node distance=1.4cm,
    input/.style={circle, draw=gray!60, fill=gray!15, minimum size=0.8cm, font=\small},
    neuron/.style={circle, draw=blue!60, fill=blue!10, minimum size=1.0cm, font=\small\bfseries},
    output/.style={circle, draw=green!60!black, fill=green!10, minimum size=0.8cm, font=\small},
    arrow/.style={-{Stealth}, thick}
]

% Entradas
\node[input] (x1) at (0, 2)   {$x_1$};
\node[input] (x2) at (0, 0.7) {$x_2$};
\node        (xd) at (0,-0.1) {$\vdots$};
\node[input] (xn) at (0,-1)   {$x_n$};

% Neurona
\node[neuron] (n) at (3, 0.7) {$\Sigma$};

% Salida
\node[output] (y) at (5.5, 0.7) {$y$};

% Bias
\node[input] (b) at (1.5, -1.8) {$b$};

% Función de activación (caja)
\node[draw=orange!70, fill=orange!10, rounded corners, minimum width=1.1cm,
      minimum height=0.7cm, font=\small] (phi) at (4.4, 0.7) {$\varphi$};

% Flechas entradas → neurona
\draw[arrow] (x1) -- node[above, font=\tiny] {$w_1$} (n);
\draw[arrow] (x2) -- node[above, font=\tiny] {$w_2$} (n);
\draw[arrow] (xn) -- node[below, font=\tiny] {$w_n$} (n);
\draw[arrow] (b)  -- (n);

% Neurona → activación → salida
\draw[arrow] (n)   -- (phi);
\draw[arrow] (phi) -- (y);

% Etiquetas
\node[font=\footnotesize, gray] at (0, -2.2) {Entradas};
\node[font=\footnotesize, gray] at (3, -2.2) {Suma ponderada};
\node[font=\footnotesize, gray] at (5.5, -2.2) {Salida};

\end{tikzpicture}
\caption{Estructura del perceptrón. Las entradas $x_i$ se combinan linealmente
con pesos $w_i$ y un sesgo $b$; la suma se pasa por una función de activación
$\varphi$ para producir la salida $y$.}
\label{fig:perceptron}
\end{figure}
